\iftestbuild % TEST-ONLY-BEGIN
\setcounter{section}{1}\setcounter{theorem}{18}
\fi % TEST-ONLY-END
% Source: book pp. 143--153 (PDF 157--167); solutions pp. 152--167.
% \axeqtag: Axler's numbered displays (5.23, 5.25, 5.30) take their number
% from the item counter; use as \[ ... \axeqtag\label{ax:5.23} \]
\DeclareRobustCommand\axeqtag{\refstepcounter{theorem}%
  \begingroup\edef\x{\endgroup\noexpand\tag{\thetheorem}}\x}

\section{The Minimal Polynomial}

\subsection{Existence of Eigenvalues on Complex Vector Spaces}

Now we come to one of the central results about operators on finite-dimensional
complex vector spaces.

\begin{result}[Existence of eigenvalues]\label{ax:5.19}
Every operator on a finite-dimensional nonzero complex vector space has an
eigenvalue.
\end{result}

\begin{proof}
Suppose $V$ is a finite-dimensional complex vector space of dimension $n>0$
and $T\in\Lin(V)$. Choose $v\in V$ with $v\ne0$. Then
\[ v,Tv,T^2v,\dots,T^nv \]
is not linearly independent, because $V$ has dimension $n$ and this list has
length $n+1$. Hence some linear combination (with not all the coefficients
equal to $0$) of the vectors above equals $0$. Thus there exists a nonconstant
polynomial $p$ of smallest degree such that
\[ p(T)v=0. \]

By the first version of the fundamental theorem of algebra (see
\ref{ax:4.12}), there exists $\lambda\in\C$ such that $p(\lambda)=0$. Hence
there exists a polynomial $q\in\Poly(\C)$ such that
\[ p(z)=(z-\lambda)q(z) \]
for every $z\in\C$ (see \ref{ax:4.6}). This implies (using \ref{ax:5.17}) that
\[ 0=p(T)v=(T-\lambda I)\bigl(q(T)v\bigr). \]
Because $q$ has smaller degree than $p$, we know that $q(T)v\ne0$. Thus the
equation above implies that $\lambda$ is an eigenvalue of $T$ with eigenvector
$q(T)v$.
\end{proof}

The proof above makes crucial use of the fundamental theorem of algebra. The
comment following Exercise~16 helps explain why the fundamental theorem of
algebra is so tightly connected to the result above.

The hypothesis in the result above that $\F=\C$ cannot be replaced with the
hypothesis that $\F=\R$, as shown by Example~\ref{ax:5.9}. The next example
shows that the finite-dimensional hypothesis in the result above also cannot
be deleted.

\begin{example}[An operator on a complex vector space with no
eigenvalues]\label{ax:5.20}
Define $T\in\Lin(\Poly(\C))$ by $(Tp)(z)=zp(z)$. If $p\in\Poly(\C)$ is a
nonzero polynomial, then the degree of $Tp$ is one more than the degree of
$p$, and thus $Tp$ cannot equal a scalar multiple of $p$. Hence $T$ has no
eigenvalues.

Because $\Poly(\C)$ is infinite-dimensional, this example does not contradict
the result above.
\end{example}

\subsection{Eigenvalues and the Minimal Polynomial}

In this subsection we introduce an important polynomial associated with each
operator. We begin with the following definition.

\begin{definition}[Monic polynomial]\label{ax:5.21}
A \emph{monic polynomial} is a polynomial whose highest-degree coefficient
equals $1$.
\end{definition}

For example, the polynomial $2+9z^2+z^7$ is a monic polynomial of degree $7$.

\begin{result}[Existence, uniqueness, and degree of minimal
polynomial]\label{ax:5.22}
Suppose $V$ is finite-dimensional and $T\in\Lin(V)$. Then there is a unique
monic polynomial $p\in\Poly(\F)$ of smallest degree such that $p(T)=0$.
Furthermore, $\deg p\le\dim V$.
\end{result}

\begin{proof}
If $\dim V=0$, then $I$ is the zero operator on $V$ and thus we take $p$ to be
the constant polynomial $1$.

Now use induction on $\dim V$. Thus assume that $\dim V>0$ and that the
desired result is true for all operators on all vector spaces of smaller
dimension. Let $u\in V$ be such that $u\ne0$. The list
$u,Tu,\dots,T^{\dim V}u$ has length $1+\dim V$ and thus is linearly dependent.
By the linear dependence lemma (\ref{ax:2.19}), there is a smallest positive
integer $m\le\dim V$ such that $T^mu$ is a linear combination of
$u,Tu,\dots,T^{m-1}u$. Thus there exist scalars
$c_0,c_1,c_2,\dots,c_{m-1}\in\F$ such that
\[ c_0u+c_1Tu+\cdots+c_{m-1}T^{m-1}u+T^mu=0. \axeqtag\label{ax:5.23} \]
Define a monic polynomial $q\in\Poly_m(\F)$ by
\[ q(z)=c_0+c_1z+\cdots+c_{m-1}z^{m-1}+z^m. \]
Then \ref{ax:5.23} implies that $q(T)u=0$.

If $k$ is a nonnegative integer, then
\[ q(T)(T^ku)=T^k\bigl(q(T)u\bigr)=T^k(0)=0. \]
The linear dependence lemma (\ref{ax:2.19}) shows that $u,Tu,\dots,T^{m-1}u$
is linearly independent. Thus the equation above implies that
$\dim\nul q(T)\ge m$. Hence
\[ \dim\range q(T)=\dim V-\dim\nul q(T)\le\dim V-m. \]
Because $\range q(T)$ is invariant under $T$ (by \ref{ax:5.18}), we can apply
our induction hypothesis to the operator $T|_{\range q(T)}$ on the vector
space $\range q(T)$. Thus there is a monic polynomial $s\in\Poly(\F)$ with
\[ \deg s\le\dim V-m \quad\text{and}\quad s(T|_{\range q(T)})=0. \]
Hence for all $v\in V$ we have
\[ \bigl((sq)(T)\bigr)(v)=s(T)\bigl(q(T)v\bigr)=0 \]
because $q(T)v\in\range q(T)$ and
$s(T)|_{\range q(T)}=s(T|_{\range q(T)})=0$. Thus $sq$ is a monic polynomial
such that $\deg sq\le\dim V$ and $(sq)(T)=0$.

The paragraph above shows that there is a monic polynomial of degree at most
$\dim V$ that when applied to $T$ gives the $0$ operator. Thus there is a monic
polynomial of smallest degree with this property, completing the existence
part of this result.

Let $p\in\Poly(\F)$ be a monic polynomial of smallest degree such that
$p(T)=0$. To prove the uniqueness part of the result, suppose $r\in\Poly(\F)$
is a monic polynomial of the same degree as $p$ and $r(T)=0$. Then
$(p-r)(T)=0$ and also $\deg(p-r)<\deg p$. If $p-r$ were not equal to $0$,
then we could divide $p-r$ by the coefficient of the highest-order term in
$p-r$ to get a monic polynomial (of smaller degree than $p$) that when applied
to $T$ gives the $0$ operator. Thus $p-r=0$, as desired.
\end{proof}

The previous result justifies the following definition.

\begin{definition}[Minimal polynomial]\label{ax:5.24}
Suppose $V$ is finite-dimensional and $T\in\Lin(V)$. Then the \emph{minimal
polynomial} of $T$ is the unique monic polynomial $p\in\Poly(\F)$ of smallest
degree such that $p(T)=0$.
\end{definition}

To compute the minimal polynomial of an operator $T\in\Lin(V)$, we need to
find the smallest positive integer $m$ such that the equation
\[ c_0I+c_1T+\cdots+c_{m-1}T^{m-1}=-T^m \]
has a solution $c_0,c_1,\dots,c_{m-1}\in\F$. If we pick a basis of $V$ and
replace $T$ in the equation above with the matrix of $T$, then the equation
above can be thought of as a system of $(\dim V)^2$ linear equations in the
$m$ unknowns $c_0,c_1,\dots,c_{m-1}\in\F$. Gaussian elimination or another
fast method of solving systems of linear equations can tell us whether a
solution exists, testing successive values $m=1,2,\dots$ until a solution
exists. By \ref{ax:5.22}, a solution exists for some smallest positive integer
$m\le\dim V$. The minimal polynomial of $T$ is then
$c_0+c_1z+\cdots+c_{m-1}z^{m-1}+z^m$.

Even faster (usually), pick $v\in V$ with $v\ne0$ and consider the equation
\[ c_0v+c_1Tv+\cdots+c_{\dim V-1}T^{\dim V-1}v=-T^{\dim V}v.
   \axeqtag\label{ax:5.25} \]
Use a basis of $V$ to convert the equation above to a system of $\dim V$
linear equations in $\dim V$ unknowns $c_0,c_1,\dots,c_{\dim V-1}$. If this
system of equations has a unique solution $c_0,c_1,\dots,c_{\dim V-1}$ (as
happens most of the time), then the scalars $c_0,c_1,\dots,c_{\dim V-1},1$ are
the coefficients of the minimal polynomial of $T$ (because \ref{ax:5.22}
states that the degree of the minimal polynomial is at most $\dim V$).

\marginnote{These estimates are based on testing millions of random
matrices.}%
Consider operators on $\R^4$ (thought of as $4$-by-$4$ matrices with respect
to the standard basis), and take $v=(1,0,0,0)$ in the paragraph above. The
faster method described above works on over $99.8\%$ of the $4$-by-$4$
matrices with integer entries in the interval $[-10,10]$ and on over
$99.999\%$ of the $4$-by-$4$ matrices with integer entries in $[-100,100]$.

The next example illustrates the faster procedure discussed above.

\begin{example}[Minimal polynomial of an operator on $\F^5$]\label{ax:5.26}
Suppose $T\in\Lin(\F^5)$ and
\[ \Mat(T)=\begin{pmatrix}
  0 & 0 & 0 & 0 & -3\\
  1 & 0 & 0 & 0 & 6\\
  0 & 1 & 0 & 0 & 0\\
  0 & 0 & 1 & 0 & 0\\
  0 & 0 & 0 & 1 & 0
\end{pmatrix} \]
with respect to the standard basis $e_1,e_2,e_3,e_4,e_5$. Taking $v=e_1$ for
\ref{ax:5.25}, we have
\begin{align*}
Te_1 &= e_2, & T^4e_1 &= T(T^3e_1)=Te_4=e_5,\\
T^2e_1 &= T(Te_1)=Te_2=e_3, & T^5e_1 &= T(T^4e_1)=Te_5=-3e_1+6e_2.\\
T^3e_1 &= T(T^2e_1)=Te_3=e_4,
\end{align*}
Thus $3e_1-6Te_1=-T^5e_1$. The list $e_1,Te_1,T^2e_1,T^3e_1,T^4e_1$, which
equals the list $e_1,e_2,e_3,e_4,e_5$, is linearly independent, so no other
linear combination of this list equals $-T^5e_1$. Hence the minimal polynomial
of $T$ is $3-6z+z^5$.
\end{example}

Recall that by definition, eigenvalues of operators on $V$ and zeros of
polynomials in $\Poly(\F)$ must be elements of $\F$. In particular, if
$\F=\R$, then eigenvalues and zeros must be real numbers.

\begin{result}[Eigenvalues are the zeros of the minimal
polynomial]\label{ax:5.27}
Suppose $V$ is finite-dimensional and $T\in\Lin(V)$.
\begin{parts}
\item The zeros of the minimal polynomial of $T$ are the eigenvalues of $T$.
\item If $V$ is a complex vector space, then the minimal polynomial of $T$ has
  the form
  \[ (z-\lambda_1)\cdots(z-\lambda_m), \]
  where $\lambda_1,\dots,\lambda_m$ is a list of all eigenvalues of $T$,
  possibly with repetitions.
\end{parts}
\end{result}

\begin{proof}
Let $p$ be the minimal polynomial of $T$.
\begin{parts}
\item First suppose $\lambda\in\F$ is a zero of $p$. Then $p$ can be written
  in the form
  \[ p(z)=(z-\lambda)q(z), \]
  where $q$ is a monic polynomial with coefficients in $\F$ (see
  \ref{ax:4.6}). Because $p(T)=0$, we have
  \[ 0=(T-\lambda I)\bigl(q(T)v\bigr) \]
  for all $v\in V$. Because $\deg q=(\deg p)-1$ and $p$ is the minimal
  polynomial of $T$, there exists at least one vector $v\in V$ such that
  $q(T)v\ne0$. The equation above thus implies that $\lambda$ is an eigenvalue
  of $T$, as desired.

  To prove that every eigenvalue of $T$ is a zero of $p$, now suppose
  $\lambda\in\F$ is an eigenvalue of $T$. Thus there exists $v\in V$ with
  $v\ne0$ such that $Tv=\lambda v$. Repeated applications of $T$ to both sides
  of this equation show that $T^kv=\lambda^kv$ for every nonnegative integer
  $k$. Thus
  \[ p(T)v=p(\lambda)v. \]
  Because $p$ is the minimal polynomial of $T$, we have $p(T)v=0$. Hence the
  equation above implies that $p(\lambda)=0$. Thus $\lambda$ is a zero of
  $p$, as desired.
\item To get the desired result, use (a) and the second version of the
  fundamental theorem of algebra (see \ref{ax:4.13}). \qedhere
\end{parts}
\end{proof}

A nonzero polynomial has at most as many distinct zeros as its degree (see
\ref{ax:4.8}). Thus (a) of the previous result, along with the result that the
minimal polynomial of an operator on $V$ has degree at most $\dim V$, gives an
alternative proof of \ref{ax:5.12}, which states that an operator on $V$ has
at most $\dim V$ distinct eigenvalues.

Every monic polynomial is the minimal polynomial of some operator, as shown by
Exercise~16, which generalizes Example~\ref{ax:5.26}. Thus \ref{ax:5.27}(a)
shows that finding exact expressions for the eigenvalues of an operator is
equivalent to the problem of finding exact expressions for the zeros of a
polynomial (and thus is not possible for some operators).

\begin{example}[An operator whose eigenvalues cannot be found
exactly]\label{ax:5.28}
Let $T\in\Lin(\C^5)$ be the operator defined by
\[ T(z_1,z_2,z_3,z_4,z_5)=(-3z_5,\,z_1+6z_5,\,z_2,\,z_3,\,z_4). \]
The matrix of $T$ with respect to the standard basis of $\C^5$ is the
$5$-by-$5$ matrix in Example~\ref{ax:5.26}. As we showed in that example, the
minimal polynomial of $T$ is the polynomial
\[ 3-6z+z^5. \]

No zero of the polynomial above can be expressed using rational numbers, roots
of rational numbers, and the usual rules of arithmetic (a proof of this would
take us considerably beyond linear algebra). Because the zeros of the
polynomial above are the eigenvalues of $T$ [by \ref{ax:5.27}(a)], we cannot
find an exact expression for any eigenvalue of $T$ in any familiar form.

Numerical techniques, which we will not discuss here, show that the zeros of
the polynomial above, and thus the eigenvalues of $T$, are approximately the
following five complex numbers:
\[ -1.67,\quad 0.51,\quad 1.40,\quad -0.12+1.59i,\quad -0.12-1.59i. \]
Note that the two nonreal zeros of this polynomial are complex conjugates of
each other, as we expect for a polynomial with real coefficients (see
\ref{ax:4.14}).
\end{example}

The next result completely characterizes the polynomials that when applied to
an operator give the $0$ operator.

\begin{result}[$q(T)=0\iff q$ is a polynomial multiple of the minimal
polynomial]\label{ax:5.29}
Suppose $V$ is finite-dimensional, $T\in\Lin(V)$, and $q\in\Poly(\F)$. Then
$q(T)=0$ if and only if $q$ is a polynomial multiple of the minimal polynomial
of $T$.
\end{result}

\begin{proof}
Let $p$ denote the minimal polynomial of $T$.

First suppose $q(T)=0$. By the division algorithm for polynomials
(\ref{ax:4.9}), there exist polynomials $s,r\in\Poly(\F)$ such that
\[ q=ps+r \axeqtag\label{ax:5.30} \]
and $\deg r<\deg p$. We have
\[ 0=q(T)=p(T)s(T)+r(T)=r(T). \]
The equation above implies that $r=0$ (otherwise, dividing $r$ by its
highest-degree coefficient would produce a monic polynomial that when applied
to $T$ gives $0$; this polynomial would have a smaller degree than the minimal
polynomial, which would be a contradiction). Thus \ref{ax:5.30} becomes the
equation $q=ps$. Hence $q$ is a polynomial multiple of $p$, as desired.

To prove the other direction, now suppose $q$ is a polynomial multiple of $p$.
Thus there exists a polynomial $s\in\Poly(\F)$ such that $q=ps$. We have
\[ q(T)=p(T)s(T)=0\,s(T)=0, \]
as desired.
\end{proof}

The next result is a nice consequence of the result above.

\begin{result}[Minimal polynomial of a restriction operator]\label{ax:5.31}
Suppose $V$ is finite-dimensional, $T\in\Lin(V)$, and $U$ is a subspace of
$V$ that is invariant under $T$. Then the minimal polynomial of $T$ is a
polynomial multiple of the minimal polynomial of $T|_U$.
\end{result}

\begin{proof}
Suppose $p$ is the minimal polynomial of $T$. Thus $p(T)v=0$ for all $v\in V$.
In particular,
\[ p(T)u=0 \text{ for all } u\in U. \]
Thus $p(T|_U)=0$. Now \ref{ax:5.29}, applied to the operator $T|_U$ in place
of $T$, implies that $p$ is a polynomial multiple of the minimal polynomial of
$T|_U$.
\end{proof}

See Exercise~25 for a result about quotient operators that is analogous to the
result above.

The next result shows that the constant term of the minimal polynomial of an
operator determines whether the operator is invertible.

\begin{result}[$T$ not invertible $\iff$ constant term of minimal polynomial
of $T$ is $0$]\label{ax:5.32}
Suppose $V$ is finite-dimensional and $T\in\Lin(V)$. Then $T$ is not
invertible if and only if the constant term of the minimal polynomial of $T$
is $0$.
\end{result}

\begin{proof}
Suppose $T\in\Lin(V)$ and $p$ is the minimal polynomial of $T$. Then
\begin{align*}
T \text{ is not invertible}
  &\iff 0 \text{ is an eigenvalue of } T\\
  &\iff 0 \text{ is a zero of } p\\
  &\iff \text{the constant term of } p \text{ is } 0,
\end{align*}
where the first equivalence holds by \ref{ax:5.7}, the second equivalence
holds by \ref{ax:5.27}(a), and the last equivalence holds because the constant
term of $p$ equals $p(0)$.
\end{proof}

\subsection{Eigenvalues on Odd-Dimensional Real Vector Spaces}

The next result will be the key tool that we use to show that every operator
on an odd-dimensional real vector space has an eigenvalue.

\begin{result}[Even-dimensional null space]\label{ax:5.33}
Suppose $\F=\R$ and $V$ is finite-dimensional. Suppose also that
$T\in\Lin(V)$ and $b,c\in\R$ with $b^2<4c$. Then $\dim\nul(T^2+bT+cI)$ is an
even number.
\end{result}

\begin{proof}
Recall that $\nul(T^2+bT+cI)$ is invariant under $T$ (by \ref{ax:5.18}). By
replacing $V$ with $\nul(T^2+bT+cI)$ and replacing $T$ with $T$ restricted to
$\nul(T^2+bT+cI)$, we can assume that $T^2+bT+cI=0$; we now need to prove
that $\dim V$ is even.

Suppose $\lambda\in\R$ and $v\in V$ are such that $Tv=\lambda v$. Then
\[ 0=(T^2+bT+cI)v=(\lambda^2+b\lambda+c)v
    =\Bigl((\lambda+\tfrac{b}{2})^2+c-\tfrac{b^2}{4}\Bigr)v. \]
The term in large parentheses above is a positive number. Thus the equation
above implies that $v=0$. Hence we have shown that $T$ has no eigenvectors.

Let $U$ be a subspace of $V$ that is invariant under $T$ and has the largest
dimension among all subspaces of $V$ that are invariant under $T$ and have
even dimension. If $U=V$, then we are done; otherwise assume there exists
$w\in V$ such that $w\notin U$.

Let $W=\spn(w,Tw)$. Then $W$ is invariant under $T$ because
$T(Tw)=-bTw-cw$. Furthermore, $\dim W=2$ because otherwise $w$ would be an
eigenvector of $T$. Now
\[ \dim(U+W)=\dim U+\dim W-\dim(U\cap W)=\dim U+2, \]
where $U\cap W=\{0\}$ because otherwise $U\cap W$ would be a one-dimensional
subspace of $V$ that is invariant under $T$ (impossible because $T$ has no
eigenvectors).

Because $U+W$ is invariant under $T$, the equation above shows that there
exists a subspace of $V$ invariant under $T$ of even dimension larger than
$\dim U$. Thus the assumption that $U\ne V$ was incorrect. Hence $V$ has even
dimension.
\end{proof}

The next result states that on odd-dimensional vector spaces, every operator
has an eigenvalue. We already know this result for finite-dimensional complex
vector spaces (without the odd hypothesis). Thus in the proof below, we will
assume that $\F=\R$.

\begin{result}[Operators on odd-dimensional vector spaces have
eigenvalues]\label{ax:5.34}
Every operator on an odd-dimensional vector space has an eigenvalue.
\end{result}

\begin{proof}
Suppose $\F=\R$ and $V$ is finite-dimensional. Let $n=\dim V$, and suppose
$n$ is an odd number. Let $T\in\Lin(V)$. We will use induction on $n$ in steps
of size two to show that $T$ has an eigenvalue. To get started, note that the
desired result holds if $\dim V=1$ because then every nonzero vector in $V$ is
an eigenvector of $T$.

Now suppose that $n\ge3$ and the desired result holds for all operators on all
odd-dimensional vector spaces of dimension less than $n$. Let $p$ denote the
minimal polynomial of $T$. If $p$ is a polynomial multiple of $x-\lambda$ for
some $\lambda\in\R$, then $\lambda$ is an eigenvalue of $T$ [by
\ref{ax:5.27}(a)] and we are done. Thus we can assume that there exist
$b,c\in\R$ such that $b^2<4c$ and $p$ is a polynomial multiple of $x^2+bx+c$
(see \ref{ax:4.16}).

There exists a monic polynomial $q\in\Poly(\R)$ such that
$p(x)=q(x)(x^2+bx+c)$ for all $x\in\R$. Now
\[ 0=p(T)=\bigl(q(T)\bigr)(T^2+bT+cI), \]
which means that $q(T)$ equals $0$ on $\range(T^2+bT+cI)$. Because
$\deg q<\deg p$ and $p$ is the minimal polynomial of $T$, this implies that
$\range(T^2+bT+cI)\ne V$.

The fundamental theorem of linear maps (\ref{ax:3.21}) tells us that
\[ \dim V=\dim\nul(T^2+bT+cI)+\dim\range(T^2+bT+cI). \]
Because $\dim V$ is odd (by hypothesis) and $\dim\nul(T^2+bT+cI)$ is even (by
\ref{ax:5.33}), the equation above shows that $\dim\range(T^2+bT+cI)$ is odd.

Hence $\range(T^2+bT+cI)$ is a subspace of $V$ that is invariant under $T$
(by \ref{ax:5.18}) and that has odd dimension less than $\dim V$. Our
induction hypothesis now implies that $T$ restricted to $\range(T^2+bT+cI)$
has an eigenvalue, which means that $T$ has an eigenvalue.
\end{proof}

See Exercise~23 in Section~8B and Exercise~10 in Section~9C for alternative
proofs of the result above.

% ------------------------------------------------------------------ exercises
\begin{exc}

\begin{exercise}
Suppose $T\in\Lin(V)$. Prove that $9$ is an eigenvalue of $T^2$ if and only if
$3$ or $-3$ is an eigenvalue of $T$.
\end{exercise}
\begin{solution}
Suppose $9$ is an eigenvalue of $T^2$. Then there exists a nonzero vector
$v\in V$ such that $T^2v=9v$. This implies $(T-3I)(T+3I)v=0$. Let
$w=(T+3I)v$, and suppose $w\ne0$. Then
\[ (T-3I)w=Tw-3w=0 \implies Tw=3w, \]
showing that $3$ is an eigenvalue of $T$. If $w=0$, then $(T+3I)v=0$, which
implies $Tv=-3v$, showing that $-3$ is an eigenvalue of $T$.

Conversely, suppose $3$ is an eigenvalue of $T$. Then there exists a nonzero
vector $v\in V$ such that $Tv=3v$. Applying $T$ again, we get
$T^2v=T(3v)=3Tv=3\cdot3v=9v$, showing that $9$ is an eigenvalue of $T^2$.
Similarly, if $-3$ is an eigenvalue of $T$, then $9$ is an eigenvalue of
$T^2$.
\end{solution}

\begin{exercise}
Suppose $V$ is a complex vector space and $T\in\Lin(V)$ has no eigenvalues.
Prove that every subspace of $V$ invariant under $T$ is either $\{0\}$ or
infinite-dimensional.
\end{exercise}
\begin{solution}
Suppose $U$ is a subspace of $V$ invariant under $T$. If $U=\{0\}$, we are
done. Suppose $U$ is nonzero and finite-dimensional. By invariance of $U$
under $T$, $T|_U\in\Lin(U)$, and $T|_U$ has an eigenvalue $\lambda\in\C$ with
corresponding eigenvector $v\in U$, $v\ne0$. But then $v$ is also an
eigenvector of $T$ with eigenvalue $\lambda$, contradicting the assumption
that $T$ has no eigenvalues. Hence, any nonzero invariant subspace $U$ must be
infinite-dimensional.
\end{solution}

\begin{exercise}
Suppose $n$ is an integer with $n>1$ and $T\in\Lin(\F^n)$ is defined by
\[ T(x_1,\dots,x_n)=(x_1+\cdots+x_n,\dots,x_1+\cdots+x_n). \]
\begin{parts}
\item Find all eigenvalues and eigenvectors of $T$.
\item Find the minimal polynomial of $T$.
\end{parts}
\exnote{The matrix of $T$ with respect to the standard basis of $\F^n$
consists of all $1$'s.}
\end{exercise}
\begin{solution}\leavevmode
\begin{parts}
\item Suppose $T(x_1,\dots,x_n)=\lambda(x_1,\dots,x_n)$ for some
  $\lambda\in\F$ and $(x_1,\dots,x_n)\ne0$. Then
  \[ \lambda x_j=x_1+\cdots+x_n \quad\text{for all } j=1,\dots,n. \]
  This implies that $\lambda x_1=\lambda x_2=\cdots=\lambda x_n$. If
  $\lambda\ne0$, then $x_1=x_2=\cdots=x_n$, and we can write
  $(x_1,\dots,x_n)=(x_1,\dots,x_1)=x_1(1,\dots,1)$. Substituting this into the
  eigenvalue equation gives
  \[ \lambda(x_1,\dots,x_1)=T(x_1,\dots,x_1)=(nx_1,\dots,nx_1), \]
  which implies $\lambda=n$. If $\lambda=0$, then $x_1+\cdots+x_n=0$, and the
  eigenvectors are all nonzero vectors $(x_1,\dots,x_n)$ such that
  $x_1+\cdots+x_n=0$. Hence, the eigenvalues of $T$ are $0$ and $n$, with
  corresponding eigenvectors as described.
\item The minimal polynomial of $T$ must have roots at the eigenvalues of $T$,
  which are $0$ and $n$. Therefore, the minimal polynomial is of the form
  $p(z)=z(z-n)$. We claim that $p(T)=0$. To verify this, let
  $x=x_1+\cdots+x_n$. Then
  \[ T(x_1,\dots,x_n)=(x,\dots,x), \]
  and
  \[ (T-nI)(x_1,\dots,x_n)=(x-nx_1,\dots,x-nx_n). \]
  Observe that $(x-nx_1)+\cdots+(x-nx_n)=nx-n(x_1+\cdots+x_n)=0$. Therefore,
  $(T-nI)(x_1,\dots,x_n)$ lies in $\nul T$, and applying $T$ again gives
  \[ T(T-nI)(x_1,\dots,x_n)=0. \]
  Hence, $p(T)=T(T-nI)=0$. To see that this is minimal, note that
  $T(1,\dots,1)=n(1,\dots,1)\ne0$, so $T\ne0$. Also,
  $(T-nI)(1,0,\dots,0)=(1-n,1,\dots,1)\ne0$, so $T-nI\ne0$. Therefore, the
  minimal polynomial of $T$ is indeed $p(z)=z(z-n)$. \qedhere
\end{parts}
\end{solution}

\begin{exercise}
Suppose $\F=\C$, $T\in\Lin(V)$, $p\in\Poly(\C)$ is a nonconstant polynomial,
and $\alpha\in\C$. Prove that $\alpha$ is an eigenvalue of $p(T)$ if and only
if $\alpha=p(\lambda)$ for some eigenvalue $\lambda$ of $T$.
\end{exercise}
\begin{solution}
Suppose $\alpha$ is an eigenvalue of $p(T)$. Then there exists a nonzero
vector $v\in V$ such that $p(T)v=\alpha v$. Then $p(T)-\alpha I$ is not
injective. By the Fundamental Theorem of Algebra, we can factor $p(z)-\alpha$
as
\[ p(z)-\alpha=c(z-\lambda_1)(z-\lambda_2)\cdots(z-\lambda_m) \]
for some $c\in\C\setminus\{0\}$ and $\lambda_1,\dots,\lambda_m\in\C$. Then
\[ p(T)-\alpha I=c(T-\lambda_1I)(T-\lambda_2I)\cdots(T-\lambda_mI). \]
Since $p(T)-\alpha I$ is not injective, at least one of the factors on the
right-hand side is not injective. Therefore, $T-\lambda_jI$ is not injective
for some $j$, which means $\lambda_j$ is an eigenvalue of $T$. Then
$p(\lambda_j)-\alpha=0$, so $\alpha=p(\lambda_j)$ for some eigenvalue
$\lambda_j$ of $T$.

Conversely, if $\alpha=p(\lambda)$ for some eigenvalue $\lambda$ of $T$ with
corresponding eigenvector $v$, then by $T^mv=\lambda^mv$ for any polynomial
$p(z)=a_mz^m+\cdots+a_1z+a_0$, we have
\[ p(T)v=a_mT^mv+\cdots+a_1Tv+a_0v=a_m\lambda^mv+\cdots+a_1\lambda v+a_0v
   =p(\lambda)v=\alpha v. \]
Hence, $\alpha$ is an eigenvalue of $p(T)$.
\end{solution}

\begin{exercise}
Give an example of an operator on $\R^2$ that shows the result in Exercise~4
does not hold if $\C$ is replaced with $\R$.
\end{exercise}
\begin{solution}
Consider the operator $T\in\Lin(\R^2)$ defined by $T(x,y)=(-y,x)$, which
represents a counterclockwise rotation by $90^\circ$. Observe that $T$ has no
real eigenvalues as follows: If $(x,y)$ is an eigenvector of $T$ with
eigenvalue $\lambda\in\R$, then $T(x,y)=\lambda(x,y)$ implies
$(-y,x)=(\lambda x,\lambda y)$. This gives the system of equations:
\[ -y=\lambda x \quad\text{and}\quad x=\lambda y. \]
Since $x$ and $y$ cannot both be zero, one can work out that $\lambda^2=-1$,
which has no real solutions. Therefore, $T$ has no real eigenvalues. Now
consider the polynomial $p(z)=z^2+1$. Then
\[ p(T)(x,y)=T^2(x,y)+(x,y)=(-x,-y)+(x,y)=(0,0). \]
Thus, $0$ is an eigenvalue of $p(T)$ with corresponding eigenvector $(1,0)$
(or any nonzero vector in $\R^2$). However, there is no real eigenvalue
$\lambda$ of $T$ such that $p(\lambda)=0$, since the eigenvalues of $T$ are
complex. This example demonstrates that the result in Exercise~5B4 does not
hold when $\C$ is replaced with $\R$.
\end{solution}

\begin{exercise}
Suppose $T\in\Lin(\F^2)$ is defined by $T(w,z)=(-z,w)$. Find the minimal
polynomial of $T$.
\end{exercise}
\begin{solution}
By the computation done in Exercise~5B5, we have $T^2(w,z)=(-w,-z)$, so
$T^2+I=0$. Therefore, the polynomial $p(z)=z^2+1$ annihilates $T$. To see that
this is the minimal polynomial, suppose $T=\alpha I$. Then
$T(1,0)=(0,1)\ne\alpha(1,0)$ for any $\alpha\in\F$, so $T\ne\alpha I$. Hence,
the minimal polynomial of $T$ is $p(z)=z^2+1$.
\end{solution}

\begin{exercise}\leavevmode
\begin{parts}
\item Give an example of $S,T\in\Lin(\F^2)$ such that the minimal polynomial
  of $ST$ does not equal the minimal polynomial of $TS$.
\item Suppose $V$ is finite-dimensional and $S,T\in\Lin(V)$. Prove that if at
  least one of $S,T$ is invertible, then the minimal polynomial of $ST$ equals
  the minimal polynomial of $TS$.
\end{parts}
\exnote{Hint: Show that if $S$ is invertible and $p\in\Poly(\F)$, then
$p(TS)=S^{-1}p(ST)S$.}
\end{exercise}
\begin{solution}\leavevmode
\begin{parts}
\item Consider the operators $S,T\in\Lin(\F^2)$ defined by $S(w,z)=(0,w)$ and
  $T(w,z)=(w,0)$. Then
  \[ ST(w,z)=S(w,0)=(0,w) \quad\text{and}\quad TS(w,z)=T(0,w)=(0,0). \]
  The minimal polynomial of $ST$ is $p(z)=z^2$, since $ST\ne0$ but
  $(ST)^2=0$, while the minimal polynomial of $TS$ is $q(z)=z$. Hence, the
  minimal polynomials of $ST$ and $TS$ are not equal.
\item Suppose $S$ is invertible. We claim that $(TS)^n=S^{-1}(ST)^nS$ for all
  positive integers $n$. We prove this by induction on $n$. For the base case
  $n=1$, we have $(TS)^1=TS=S^{-1}(ST)S$. Now suppose the claim holds for some
  positive integer $k$, i.e., $(TS)^k=S^{-1}(ST)^kS$. Then
  \[ (TS)^{k+1}=(TS)^k(TS)=S^{-1}(ST)^kS(TS)=S^{-1}(ST)^k(ST)S
     =S^{-1}(ST)^{k+1}S. \]
  By induction, the claim holds for all positive integers $n$. Let
  $p=z^m+a_{m-1}z^{m-1}+\cdots+a_1z+a_0$ be the minimal polynomial of $ST$.
  Then
  \begin{align*}
  p(TS) &= (TS)^m+a_{m-1}(TS)^{m-1}+\cdots+a_1(TS)+a_0I\\
        &= S^{-1}(ST)^mS+a_{m-1}S^{-1}(ST)^{m-1}S+\cdots+a_1S^{-1}(ST)S
           +S^{-1}(a_0I)S\\
        &= S^{-1}\bigl((ST)^m+a_{m-1}(ST)^{m-1}+\cdots+a_1(ST)+a_0I\bigr)S\\
        &= S^{-1}p(ST)S=0,
  \end{align*}
  where $S^{-1}S=I$ implies $S^{-1}(a_0I)S=a_0I$. Hence, $p$ annihilates $TS$,
  and since $p$ is the minimal polynomial of $ST$, it follows that the minimal
  polynomial $q$ of $TS$ divides $p$. Observing that $ST=S(TS)S^{-1}$, we can
  similarly show that $q$ annihilates $ST$, and thus $p$ divides $q$.
  Therefore, $p=q$, and the minimal polynomials of $ST$ and $TS$ are equal.
  \qedhere
\end{parts}
\end{solution}

\begin{exercise}
Suppose $T\in\Lin(\R^2)$ is the operator of counterclockwise rotation by
$1^\circ$. Find the minimal polynomial of $T$.
\exnote{Because $\dim\R^2=2$, the degree of the minimal polynomial of $T$ is
at most $2$. Thus the minimal polynomial of $T$ is not the tempting polynomial
$x^{180}+1$, even though $T^{180}=-I$.}
\end{exercise}
\begin{solution}
Since $T\ne0$, the minimal polynomial of $T$ is not linear. Moreover,
$\dim\R^2=2$ implies that the degree of the minimal polynomial of $T$ must be
$2$. Therefore, the minimal polynomial of $T$ is of the form
\[ p(z)=z^2+bz+c. \]
Consider the point $(1,0)\in\R^2$. Note that
$T(x,y)=(x\cos1^\circ-y\sin1^\circ,\,x\sin1^\circ+y\cos1^\circ)$. Then
\[ T(1,0)=(\cos1^\circ,\sin1^\circ),\quad T^2(1,0)=(\cos2^\circ,\sin2^\circ). \]
Since $p(T)=0$, we have
\begin{align*}
p(T)(1,0) &= T^2(1,0)+bT(1,0)+c(1,0)\\
          &= (\cos2^\circ+b\cos1^\circ+c,\,\sin2^\circ+b\sin1^\circ)=(0,0).
\end{align*}
This gives the system of equations:
\[ \begin{cases} \cos2^\circ+b\cos1^\circ+c=0,\\
   \sin2^\circ+b\sin1^\circ=0. \end{cases} \]
From this, we get the values of $b$ and $c$:
\[ b=-\frac{\sin2^\circ}{\sin1^\circ}=-2\cos1^\circ,\quad
   c=-\cos2^\circ-b\cos1^\circ=1. \]
Therefore, the minimal polynomial of $T$ is $p(z)=z^2-2\cos1^\circ z+1$.
\end{solution}

\begin{exercise}
Suppose $T\in\Lin(V)$ is such that with respect to some basis of $V$, all
entries of the matrix of $T$ are rational numbers. Explain why all
coefficients of the minimal polynomial of $T$ are rational numbers.
\end{exercise}
\begin{solution}
Let $A$ be the matrix of $T$ with respect to some basis of $V$. Since all
entries of $A$ are rational numbers, $A^k$ also has rational entries for any
positive integer $k$. Then the minimal polynomial
$p(z)=z^m+a_{m-1}z^{m-1}+\cdots+a_1z+a_0$ of $T$ in terms of the matrix $A$
satisfies
\[ p(A)=A^m+a_{m-1}A^{m-1}+\cdots+a_1A+a_0I=0. \]
Then this gives a system of $(\dim V)^2$ linear equations in the coefficients
$a_0,a_1,\allowbreak\dots,a_{m-1}$ with rational coefficients. By uniqueness of the
minimal polynomial, this system has a unique solution. Since the system has
rational coefficients and a unique solution, the solution must also be
rational. Therefore, all coefficients of the minimal polynomial of $T$ are
rational numbers.
\end{solution}

\begin{exercise}
Suppose $V$ is finite-dimensional, $T\in\Lin(V)$, and $v\in V$. Prove that
\[ \spn(v,Tv,\dots,T^mv)=\spn(v,Tv,\dots,T^{\dim V-1}v) \]
for all integers $m\ge\dim V-1$.
\end{exercise}
\begin{solution}
Let $n=\dim V$ and let $p(z)=z^d+a_{d-1}z^{d-1}+\cdots+a_0$ be the minimal
polynomial of $T$. Then $d\le n$ and $p(T)=0$. For any $v\in V$,
\begin{multline*}
p(T)v=0 \implies T^dv=-a_{d-1}T^{d-1}v-\cdots-a_1Tv-a_0v\\
  \implies T^dv\in\spn(v,Tv,\dots,T^{d-1}v).
\end{multline*}
Let $U=\spn(v,Tv,\dots,T^{d-1}v)$.

We claim that $U$ is invariant under $T$. Take any
$u=c_0v+c_1Tv+\cdots+c_{d-1}T^{d-1}v\in U$. Then,
\[ Tu=c_0Tv+c_1T^2v+\cdots+c_{d-2}T^{d-1}v+c_{d-1}T^dv. \]
Every term on the right-hand side except $T^dv$ is automatically in the
spanning set of $U$. We also showed that $T^dv\in U$ earlier, so $Tu\in U$.
Hence, $T(U)\subseteq U$.

Because $v\in U$ and $U$ is invariant under $T$, we see that $T^mv\in U$ for
every $m\ge0$ (this is immediate by induction). So, for any $m\ge n-1$, we
have that
\[ \spn(v,Tv,\dots,T^mv)\subseteq U\subseteq\spn(v,Tv,\dots,T^{n-1}v) \]
by using $d-1\le n-1$ to obtain the second inclusion. Since
$\{v,Tv,\dots,T^{n-1}v\}\subseteq\{v,Tv,\dots,T^mv\}$, the reverse inclusion
is immediate, and so
\[ \spn(v,Tv,\dots,T^mv)=\spn(v,Tv,\dots,T^{n-1}v). \qedhere \]
\end{solution}

\begin{exercise}
Suppose $V$ is a two-dimensional vector space, $T\in\Lin(V)$, and the matrix
of $T$ with respect to some basis of $V$ is
$\begin{pmatrix} a & c\\ b & d \end{pmatrix}$.
\begin{parts}
\item Show that $T^2-(a+d)T+(ad-bc)I=0$.
\item Show that the minimal polynomial of $T$ equals
  \[ \begin{cases} z-a & \text{if } b=c=0 \text{ and } a=d,\\
     z^2-(a+d)z+(ad-bc) & \text{otherwise}. \end{cases} \]
\end{parts}
\end{exercise}
\begin{solution}\leavevmode
\begin{parts}
\item We compute $T^2$:
  \[ T^2=\begin{pmatrix} a & c\\ b & d \end{pmatrix}
         \begin{pmatrix} a & c\\ b & d \end{pmatrix}
        =\begin{pmatrix} a^2+bc & ac+cd\\ ab+bd & bc+d^2 \end{pmatrix}. \]
  Then
  \begin{align*}
  &T^2-(a+d)T+(ad-bc)I\\
  &\quad=\begin{pmatrix} a^2+bc & ac+cd\\ ab+bd & bc+d^2 \end{pmatrix}
     -(a+d)\begin{pmatrix} a & c\\ b & d \end{pmatrix}
     +(ad-bc)\begin{pmatrix} 1 & 0\\ 0 & 1 \end{pmatrix}\\
  &\quad=\begin{pmatrix} a^2+bc-(a+d)a+ad-bc & ac+cd-(a+d)c\\
     ab+bd-(a+d)b & bc+d^2-(a+d)d+ad-bc \end{pmatrix}\\
  &\quad=\begin{pmatrix} a^2+bc-a^2-ad+ad-bc & ac+cd-ac-cd\\
     ab+bd-ab-bd & bc+d^2-d^2-ad+ad-bc \end{pmatrix}\\
  &\quad=\begin{pmatrix} 0 & 0\\ 0 & 0 \end{pmatrix}.
  \end{align*}
\item If $b=c=0$ and $a=d$, then the matrix of $T$ is
  \[ \begin{pmatrix} a & 0\\ 0 & a \end{pmatrix}, \]
  which is a scalar multiple of the identity matrix. In this case, the minimal
  polynomial of $T$ is $z-a$.

  Otherwise, the polynomial $z^2-(a+d)z+(ad-bc)$ annihilates $T$ by part (a).
  To see that this is the minimal polynomial, we need to check that $T$ is not
  a scalar multiple of the identity matrix. If $T$ were a scalar multiple of
  the identity matrix, then we would have $b=c=0$ and $a=d$, which contradicts
  our assumption. Therefore, the minimal polynomial of $T$ has degree $2$ and
  is $z^2-(a+d)z+(ad-bc)$ in this case. \qedhere
\end{parts}
\end{solution}

\begin{exercise}
Define $T\in\Lin(\F^n)$ by
$T(x_1,x_2,x_3,\dots,x_n)=(x_1,2x_2,3x_3,\dots,nx_n)$. Find the minimal
polynomial of $T$.
\end{exercise}
\begin{solution}
By Exercise~5A42, the eigenvalues of $T$ are $1,2,\dots,n$, whose eigenvectors
are scalar multiples of the standard basis vectors $e_1,e_2,\dots,e_n$. The
minimal polynomial of $T$ must have roots at the eigenvalues of $T$, which are
$1,2,\dots,n$. Therefore, the minimal polynomial is of the form
$p(z)=(z-1)(z-2)\cdots(z-n)$. To verify that this is indeed the minimal
polynomial, we can check that $p(T)=0$. Observe that $p(T)e_j=0$ for each
$j=1,2,\dots,n$, since $Te_j=je_j$ and thus $(T-jI)e_j=0$. Since
$\{e_1,e_2,\dots,e_n\}$ establishes a basis for $\F^n$, it follows that
$p(T)=0$. Therefore, the minimal polynomial of $T$ is divisible by
$(z-1)(z-2)\cdots(z-n)$ and has degree $n$.

To see that this is minimal, note that if we remove any factor $(z-k)$ from
$p(z)$, the resulting polynomial would not annihilate $T$ and would result in
$(k-1)(k-2)\cdots\bigl(k-(k-1)\bigr)\bigl(k-(k+1)\bigr)\cdots(k-n)\ne0$ when
applied to $e_k$. Therefore, the minimal polynomial of $T$ is
$p(z)=(z-1)(z-2)\cdots(z-n)$.
\end{solution}

\begin{exercise}
Suppose $V$ is finite-dimensional, $T\in\Lin(V)$, and $p\in\Poly(\F)$. Prove
that there exists a unique $r\in\Poly(\F)$ such that $p(T)=r(T)$ and $\deg r$
is less than the degree of the minimal polynomial of $T$.
\end{exercise}
\begin{solution}
Let $m$ be the minimal polynomial of $T$ with degree $d$. By the Division
Algorithm for polynomials, we can write
\[ p(z)=q(z)m(z)+r(z), \]
where $q,r\in\Poly(\F)$ and $\deg r<d$. Applying this to $T$, we have
\[ p(T)=q(T)m(T)+r(T). \]
Since $m(T)=0$, it follows that
\[ p(T)=r(T). \]
Thus, there exists a polynomial $r$ such that $p(T)=r(T)$ and $\deg r<d$.

To show uniqueness, suppose there exists another polynomial $s\in\Poly(\F)$
such that $p(T)=s(T)$ and $\deg s<d$. Then we have
\[ r(T)=p(T)=s(T). \]
This implies that $(r-s)(T)=0$. Since $\deg(r-s)<d$, it follows that $r-s$
must be the zero polynomial (otherwise, it would contradict the minimality of
$m$). Hence, $r=s$, proving uniqueness.
\end{solution}

\begin{exercise}
Suppose $V$ is finite-dimensional and $T\in\Lin(V)$ has minimal polynomial
$4+5z-6z^2-7z^3+2z^4+z^5$. Find the minimal polynomial of $T^{-1}$.
\end{exercise}
\begin{solution}
Let $p(z)=4+5z-6z^2-7z^3+2z^4+z^5$ be the minimal polynomial of $T$. We want
to find the minimal polynomial of $T^{-1}$. Since $p(z)$ is the minimal
polynomial of $T$, we have
\[ p(T)=4I+5T-6T^2-7T^3+2T^4+T^5=0. \]
To find the minimal polynomial of $T^{-1}$, apply $T^{-1}$ to both sides of
the equation until we express it in terms of $T^{-1}$:
\[ 4T^{-5}+5T^{-4}-6T^{-3}-7T^{-2}+2T^{-1}+I=0. \]
Then the polynomial $q(z)=4z^5+5z^4-6z^3-7z^2+2z+1$ annihilates $T^{-1}$. To
make it monic, we can divide by the leading coefficient:
\[ q(z)=z^5+5/4z^4-3/2z^3-7/4z^2+1/2z+1/4. \]
Moreover, suppose $T^{-1}$ has a minimal polynomial $r(z)$ of degree less than
$5$. Then $r(T^{-1})=0$, which implies that $r(T)=0$ as well, contradicting
the minimality of $p(z)$. Hence, the minimal polynomial of $T^{-1}$ is $q(z)$.
\end{solution}

\begin{exercise}
Suppose $V$ is a finite-dimensional complex vector space with $\dim V>0$ and
$T\in\Lin(V)$. Define $f\colon\C\to\R$ by
\[ f(\lambda)=\dim\range(T-\lambda I). \]
Prove that $f$ is not a continuous function.
\end{exercise}
\begin{solution}
Since $V$ is a finite-dimensional complex vector space, the operator $T$ has
at least one eigenvalue $\lambda_0\in\C$. By Exercise~5A11, there exists
$\delta>0$ such that $T-\lambda I$ is invertible for all $\lambda\in\C$ with
$0<\abs{\lambda-\lambda_0}<\delta$. This implies that for all such $\lambda$,
we have $\dim\range(T-\lambda I)=\dim V$. However, at the eigenvalue
$\lambda_0$, we have $\dim\range(T-\lambda_0I)<\dim V$ since the nullity is
positive. Therefore, the function $f$ has a jump discontinuity at
$\lambda_0$, and thus $f$ is not continuous.
\end{solution}

\begin{exercise}
Suppose $a_0,\dots,a_{n-1}\in\F$. Let $T$ be the operator on $\F^n$ whose
matrix (with respect to the standard basis) is
\[ \begin{pmatrix}
  0 &   &        &   & -a_0\\
  1 & 0 &        &   & -a_1\\
    & 1 & \ddots &   & -a_2\\
    &   & \ddots &   & \vdots\\
    &   &        & 0 & -a_{n-2}\\
    &   &        & 1 & -a_{n-1}
\end{pmatrix}. \]
Here all entries of the matrix are $0$ except for all $1$'s on the line under
the diagonal and the entries in the last column (some of which might also be
$0$). Show that the minimal polynomial of $T$ is the polynomial
\[ a_0+a_1z+\cdots+a_{n-1}z^{n-1}+z^n. \]
\exnote{The matrix above is called the \textbf{companion matrix} of the
polynomial above. This exercise shows that every monic polynomial is the
minimal polynomial of some operator. Hence a formula or an algorithm that
could produce exact eigenvalues for each operator on each $\F^n$ could then
produce exact zeros for each polynomial [by \ref{ax:5.27}(a)]. Thus there is
no such formula or algorithm. However, efficient numerical methods exist for
obtaining very good approximations for the eigenvalues of an operator.}
\end{exercise}
\begin{solution}
Let $p(z)=a_0+a_1z+\cdots+a_{n-1}z^{n-1}+z^n$. Let $e_1,\dots,e_n$ be the
standard basis of $\F^n$. We will show that $p(T)=0$ and that no polynomial of
lower degree annihilates $T$. First, we compute $T^ke_1$ for $k=0,1,\dots,n$:
\begin{gather*}
T^0e_1=e_1,\quad T^1e_1=e_2,\quad T^2e_1=e_3,\quad\dots,\quad
T^{n-1}e_1=e_n,\\
T^ne_1=-a_0e_1-a_1e_2-\cdots-a_{n-1}e_n.
\end{gather*}
Then we have
\begin{align*}
p(T)e_1 &= a_0T^0e_1+a_1T^1e_1+\cdots+a_{n-1}T^{n-1}e_1+T^ne_1\\
        &= a_0e_1+a_1e_2+\cdots+a_{n-1}e_n-a_0e_1-a_1e_2-\cdots-a_{n-1}e_n\\
        &= 0.
\end{align*}
Since $T^{j-1}e_1=e_j$ for $j=1,2,\dots,n$, and $p(T)$ commutes with $T$, we
have
\[ p(T)e_j=p(T)T^{j-1}e_1=T^{j-1}p(T)e_1=T^{j-1}0=0 \]
for all $j=1,2,\dots,n$. Since $\{e_1,e_2,\dots,e_n\}$ is a basis for $\F^n$,
it follows that $p(T)=0$.

To show that $p$ is the minimal polynomial, suppose there exists a polynomial
$q(z)=b_0+b_1z+\cdots+b_{m-1}z^{m-1}+z^m$ of degree $m<n$ such that $q(T)=0$.
Then we would have
\[ q(T)e_1=b_0e_1+b_1e_2+\cdots+b_{m-1}e_m+T^me_1=0. \]
However, $T^me_1=e_{m+1}$, which is not in the span of $\{e_1,e_2,\dots,e_m\}$
since $m<n$. This leads to a contradiction, as $e_{m+1}$ cannot be expressed
as a linear combination of $e_1,e_2,\dots,e_m$. Therefore, no polynomial of
degree less than $n$ annihilates $T$, and the minimal polynomial of $T$ is
indeed $p(z)=a_0+a_1z+\cdots+a_{n-1}z^{n-1}+z^n$.
\end{solution}

\begin{exercise}
Suppose $V$ is finite-dimensional, $T\in\Lin(V)$, and $p$ is the minimal
polynomial of $T$. Suppose $\lambda\in\F$. Show that the minimal polynomial of
$T-\lambda I$ is the polynomial $q$ defined by $q(z)=p(z+\lambda)$.
\end{exercise}
\begin{solution}
Let $p(z)$ be the minimal polynomial of $T$. Then we have
\[ p(T)=0. \]
Consider the polynomial $q(z)=p(z+\lambda)$. Then we have
\[ q(T-\lambda I)=p\bigl((T-\lambda I)+\lambda I\bigr)=p(T)=0. \]
Thus, $q$ annihilates $T-\lambda I$.

To show that $q$ is the minimal polynomial, suppose there exists a polynomial
$r(z)$ of degree less than $\deg q$ such that $r(T-\lambda I)=0$. Then we can
write
\[ r(z)=s(z+\lambda) \]
for some polynomial $s(z)$ of degree less than $\deg p$. Then we have
\[ r(T-\lambda I)=s\bigl((T-\lambda I)+\lambda I\bigr)=s(T)=0. \]
However, this contradicts the minimality of $p$, since $\deg s<\deg p$ and
$s(T)=0$. Therefore, no polynomial of degree less than $\deg q$ annihilates
$T-\lambda I$, and the minimal polynomial of $T-\lambda I$ is indeed
$q(z)=p(z+\lambda)$.
\end{solution}

\begin{exercise}
Suppose $V$ is finite-dimensional, $T\in\Lin(V)$, and $p$ is the minimal
polynomial of $T$. Suppose $\lambda\in\F\setminus\{0\}$. Show that the minimal
polynomial of $\lambda T$ is the polynomial $q$ defined by
$q(z)=\lambda^{\deg p}\,p\Bigl(z/\lambda\Bigr)$.
\end{exercise}
\begin{solution}
Let $p(z)$ be the minimal polynomial of $T$. Then we have
\[ p(T)=0. \]
Consider the polynomial $q(z)=\lambda^{\deg p}p(z/\lambda)$. Then we
have
\[ q(\lambda T)=\lambda^{\deg p}p\Bigl(\frac{\lambda T}{\lambda}\Bigr)
   =\lambda^{\deg p}p(T)=0. \]
Moreover, $q$ has leading term
$\lambda^{\deg p}(z/\lambda)^{\deg p}=z^{\deg p}$, so $q$ is monic and
has the same degree as $p$. To show that $q$ is the minimal polynomial,
suppose there exists a polynomial $r(z)$ of degree less than $\deg q$ such
that $r(\lambda T)=0$. Define $s(z)=r(\lambda z)$. Then we have
\[ s(T)=r(\lambda T)=0. \]
However, $\deg s=\deg r<\deg q=\deg p$, which contradicts the minimality of
$p$. Therefore, no polynomial of degree less than $\deg q$ annihilates
$\lambda T$, and the minimal polynomial of $\lambda T$ is indeed
$q(z)=\lambda^{\deg p}p(z/\lambda)$.
\end{solution}

\begin{exercise}
Suppose $V$ is finite-dimensional and $T\in\Lin(V)$. Let $\mathcal{E}$ be the
subspace of $\Lin(V)$ defined by
\[ \mathcal{E}=\{q(T) : q\in\Poly(\F)\}. \]
Prove that $\dim\mathcal{E}$ equals the degree of the minimal polynomial of
$T$.
\end{exercise}
\begin{solution}
Let $p(z)$ be the minimal polynomial of $T$ with degree $d$. We will show that
$\dim\mathcal{E}=d$.

First, we show that $\dim\mathcal{E}\le d$. By the Division Algorithm for
polynomials, for any polynomial $q\in\Poly(\F)$, we can write
\[ q(z)=s(z)p(z)+r(z), \]
where $s,r\in\Poly(\F)$ and $\deg r<d$. Then we have
\[ q(T)=s(T)p(T)+r(T)=0+r(T)=r(T). \]
Since $\deg r<d$, the set $\{I,T,T^2,\dots,T^{d-1}\}$ spans $\mathcal{E}$.
Hence, $\dim\mathcal{E}\le d$.

Next, we show that $\dim\mathcal{E}\ge d$. Suppose there exists a nontrivial
linear combination of $\{I,T,T^2,\dots,T^{d-1}\}$ that equals zero:
\[ a_0I+a_1T+a_2T^2+\cdots+a_{d-1}T^{d-1}=0. \]
This would imply that the polynomial
\[ q(z)=a_0+a_1z+a_2z^2+\cdots+a_{d-1}z^{d-1} \]
annihilates $T$, contradicting the minimality of $p(z)$ since $\deg q<d$.
Therefore, the set $\{I,T,T^2,\dots,T^{d-1}\}$ is linearly independent in
$\mathcal{E}$, which implies that $\dim\mathcal{E}\ge d$.

Combining both inequalities, we conclude that $\dim\mathcal{E}=d$, which is
the degree of the minimal polynomial of $T$.
\end{solution}

\begin{exercise}
Suppose $T\in\Lin(\F^4)$ is such that the eigenvalues of $T$ are $3,5,8$.
Prove that $(T-3I)^2(T-5I)^2(T-8I)^2=0$.
\end{exercise}
\begin{solution}
Let $p(z)$ be the minimal polynomial of $T$. Since the eigenvalues of $T$ are
$3,5,8$, the minimal polynomial must have these roots. Therefore, we can write
\[ p(z)=(z-3)^{m_1}(z-5)^{m_2}(z-8)^{m_3}, \]
where $m_1,m_2,m_3\ge1$ are the multiplicities of the eigenvalues in the
minimal polynomial.

Since $T\in\Lin(\F^4)$, the degree of the minimal polynomial cannot exceed
$\dim\F^4=4$. Thus, we have
\[ m_1+m_2+m_3\le4. \]

To show that $(T-3I)^2(T-5I)^2(T-8I)^2=0$, we need to show that each
$m_i\le2$. Without loss of generality, suppose $m_1\ge3$. Then
$m_2+m_3\le1$, which implies that at least one of $m_2$ or $m_3$ is equal to
$0$, contradicting the fact that both $5$ and $8$ are eigenvalues of $T$.
Hence, $m_1\le2$. By similar reasoning, we can conclude that $m_2\le2$ and
$m_3\le2$.

Therefore, we have $m_1,m_2,m_3\le2$. Then $p(z)$ divides
$(z-3)^2(z-5)^2(z-8)^2$ for some polynomial $q(z)$, and thus
\[ (T-3I)^2(T-5I)^2(T-8I)^2=q(T)p(T)=q(T)\cdot0=0. \qedhere \]
\end{solution}

\begin{exercise}
Suppose $V$ is finite-dimensional and $T\in\Lin(V)$. Prove that the minimal
polynomial of $T$ has degree at most $1+\dim\range T$.
\exnote{If $\dim\range T<\dim V-1$, then this exercise gives a better upper
bound than \ref{ax:5.22} for the degree of the minimal polynomial of $T$.}
\end{exercise}
\begin{solution}
Let $r=\dim\range T$. Observe that $\range T$ is invariant under $T$, and we
can consider the restriction of $T$ to $\range T$, denoted by
$T|_{\range T}$. Let $p(z)$ be the minimal polynomial of $T$ and let $q(z)$ be
the minimal polynomial of $T|_{\range T}$. Then $\deg q\le r$ since
$T|_{\range T}$ is an operator on a space of dimension $r$.

We now show that the polynomial $zq(z)$ annihilates $T$. For any $v\in V$, we
know that $Tv\in\range T$. By invariance, we know that $w\in\range T$ implies
$(T|_{\range T})^kw\in\range T$ for any $k\ge0$. Then
\[ q(T)w=q(T|_{\range T})w=0. \]
Therefore, for any $v\in V$, we have
\[ zq(z)(T)v=Tq(T)v=q(T)Tv=T\cdot0=0, \]
where we used the fact that $T$ commutes with $q(T)$ since $q(T)$ is a
polynomial in $T$. This shows that $zq(z)$ annihilates $T$, and hence the
minimal polynomial $p(z)$ of $T$ divides $zq(z)$. Since $\deg q\le r$, we have
$\deg p\le1+r=1+\dim\range T$.
\end{solution}

\begin{exercise}
Suppose $V$ is finite-dimensional and $T\in\Lin(V)$. Prove that $T$ is
invertible if and only if $I\in\spn(T,T^2,\dots,T^{\dim V})$.
\end{exercise}
\begin{solution}
Suppose $T$ is invertible. Let $q(z)$ be the minimal polynomial of $T$. Since
$T$ is invertible, $0$ is not an eigenvalue of $T$, and thus $q(0)\ne0$. We
may write $q(z)$ as
\[ q(z)=z^m+a_{m-1}z^{m-1}+\cdots+a_1z+a_0, \]
where $a_0\ne0$ and $m\le\dim V$. Then we have
\[ 0=q(T)=T^m+a_{m-1}T^{m-1}+\cdots+a_1T+a_0I. \]
Rearranging gives
\begin{align*}
I &= -\frac{1}{a_0}(T^m+a_{m-1}T^{m-1}+\cdots+a_1T)\\
  &\in \spn(T,T^2,\dots,T^m)\subseteq\spn(T,T^2,\dots,T^{\dim V}).
\end{align*}
Hence, if $T$ is invertible, then $I\in\spn(T,T^2,\dots,T^{\dim V})$.

Conversely, suppose $I\in\spn(T,T^2,\dots,T^{\dim V})$. Then there exist
scalars $a_1,a_2,\allowbreak\dots,a_{\dim V}$ such that
\[ I=a_1T+a_2T^2+\cdots+a_{\dim V}T^{\dim V}. \]
Rearranging gives
\[ 0=-I+a_1T+a_2T^2+\cdots+a_{\dim V}T^{\dim V}=p(T), \]
where $p(z)=-1+a_1z+a_2z^2+\cdots+a_{\dim V}z^{\dim V}$. Since the minimal
polynomial of $T$ divides any polynomial that annihilates $T$, the minimal
polynomial of $T$ must also have a nonzero constant term. Therefore, $0$ is
not an eigenvalue of $T$, and hence $T$ is invertible.
\end{solution}

\begin{exercise}
Suppose $V$ is finite-dimensional and $T\in\Lin(V)$. Let $n=\dim V$. Prove
that if $v\in V$, then $\spn(v,Tv,\dots,T^{n-1}v)$ is invariant under $T$.
\end{exercise}
\begin{solution}
Let $W=\spn(v,Tv,\dots,T^{n-1}v)$. To show that $W$ is invariant under $T$, we
may show that $Tv,T^2v,\dots,T^nv\in W$. Observe that $Tv$, $T(Tv)=T^2v$,
$T(T^2v)=T^3v$, \dots, $T(T^{n-2}v)=T^{n-1}v$ are all in $W$ by definition. It
remains to show that $T^nv\in W$.

Let $p(z)$ be the minimal polynomial of $T$. Let $m=\deg p$. Since $m\le n$,
we have
\[ p(T)v=T^mv+a_{m-1}T^{m-1}v+\cdots+a_1Tv+a_0v=0. \]
Rearranging gives
\[ T^mv=-a_{m-1}T^{m-1}v-\cdots-a_1Tv-a_0v\in W. \]
If $m=n$, then we are done. If $m<n$, then we can apply $T^{n-m}$ to both
sides to get
\[ T^nv=-a_{m-1}T^{n-1}v-\cdots-a_1T^{n-m+1}v-a_0T^{n-m}v\in W, \]
where we used the fact that $T^{n-m}v,T^{n-m+1}v,\dots,T^{n-1}v\in W$ since
$n-m<n$. Therefore, $T^nv\in W$, and hence $W$ is invariant under $T$.
\end{solution}

\begin{exercise}
Suppose $V$ is a finite-dimensional complex vector space. Suppose
$T\in\Lin(V)$ is such that $5$ and $6$ are eigenvalues of $T$ and that $T$ has
no other eigenvalues. Prove that $(T-5I)^{\dim V-1}(T-6I)^{\dim V-1}=0$.
\end{exercise}
\begin{solution}
Let $n=\dim V$. Since $5$ and $6$ are the only eigenvalues of $T$, the minimal
polynomial of $T$ must be of the form
\[ p(z)=(z-5)^{m_1}(z-6)^{m_2}, \]
where $m_1,m_2\ge1$ are the multiplicities of the eigenvalues in the minimal
polynomial.

Since $T\in\Lin(V)$, we have $\deg p\le n$. Therefore, we have
\[ m_1+m_2\le n. \]

To show that $(T-5I)^{n-1}(T-6I)^{n-1}=0$, we need to show that each
$m_i\le n-1$. Suppose $m_1\ge n$. Then $m_2\le0$, which contradicts the fact
that $6$ is an eigenvalue of $T$. Hence, $m_1\le n-1$. By similar reasoning,
we can conclude that $m_2\le n-1$.

Therefore, we have $m_1,m_2\le n-1$. Then $p(z)$ divides
$(z-5)^{n-1}(z-6)^{n-1}$, so there exists a polynomial $q(z)$ such that
$p(z)q(z)=(z-5)^{n-1}(z-6)^{n-1}$. Applying this to $T$, we have
\[ (T-5I)^{n-1}(T-6I)^{n-1}=p(T)q(T)=0\cdot q(T)=0. \qedhere \]
\end{solution}

\begin{exercise}
Suppose $V$ is finite-dimensional, $T\in\Lin(V)$, and $U$ is a subspace of $V$
that is invariant under $T$.
\begin{parts}
\item Prove that the minimal polynomial of $T$ is a polynomial multiple of the
  minimal polynomial of the quotient operator $T/U$.
\item Prove that
  \[ (\text{minimal polynomial of } T|_U)\times(\text{minimal polynomial of }
     T/U) \]
  is a polynomial multiple of the minimal polynomial of $T$.
\end{parts}
\exnote{The quotient operator $T/U$ was defined in Exercise~38 in
Section~5A.}
\end{exercise}
\begin{solution}\leavevmode
\begin{parts}
\item Let $p(z)$ be the minimal polynomial of $T$ and let $q(z)$ be the
  minimal polynomial of $T/U$. Recall that $T/U$ is defined by
  $(T/U)(v+U)=Tv+U$ for all $v\in V$. Inductively, we have
  $(T/U)^k(v+U)=T^kv+U$ for all $k\ge0$. Then we have
  \[ p(T/U)(v+U)=p(T)v+U=0+U=U, \]
  where we used the fact that $p(T)=0$. Therefore, $p(T/U)=0$, which implies
  that $q(z)$ divides $p(z)$. Hence, the minimal polynomial of $T$ is a
  polynomial multiple of the minimal polynomial of the quotient operator
  $T/U$.
\item Let $p(z)$ be the minimal polynomial of $T$, let $q(z)$ be the minimal
  polynomial of $T|_U$, and let $r(z)$ be the minimal polynomial of $T/U$. By
  definition,
  \[ r(T/U)(v+U)=r(T)v+U=U, \]
  where we used the fact that $r(T/U)=0$. Then $r(T)v\in U$ for all $v\in V$.
  Moreover, invariance of $U$ under $T$ implies that $q(T)u=q(T|_U)u=0$ for
  all $u\in U$. Since $r(T)v\in U$, and $q(T)=q(T|_U)$ on $U$, we have
  \[ q(T)r(T)v=q(T|_U)r(T)v=0 \]
  for all $v\in V$. Therefore, $q(T)r(T)=0$, which implies $p(z)$ divides
  $q(z)r(z)$. Hence, the product of the minimal polynomials of $T|_U$ and
  $T/U$ is a polynomial multiple of the minimal polynomial of $T$. \qedhere
\end{parts}
\end{solution}

\begin{exercise}
Suppose $V$ is finite-dimensional, $T\in\Lin(V)$, and $U$ is a subspace of $V$
that is invariant under $T$. Prove that the set of eigenvalues of $T$ equals
the union of the set of eigenvalues of $T|_U$ and the set of eigenvalues of
$T/U$.
\end{exercise}
\begin{solution}
Let $\lambda$ be an eigenvalue of $T$ with corresponding eigenvector
$v\in V\setminus\{0\}$. Then we have
\[ Tv=\lambda v. \]
If $v\in U$, then $\lambda$ is an eigenvalue of $T|_U$. If $v\notin U$, then
consider the coset $v+U\in V/U$. We have
\[ (T/U)(v+U)=Tv+U=\lambda v+U=\lambda(v+U), \]
which shows that $\lambda$ is an eigenvalue of $T/U$.

Conversely, suppose $\lambda$ is an eigenvalue of $T|_U$ with corresponding
eigenvector $u\in U\setminus\{0\}$. Then we have
\[ T|_Uu=Tu=\lambda u, \]
which shows that $\lambda$ is also an eigenvalue of $T$.

Now suppose $\lambda$ is an eigenvalue of $T/U$ with corresponding eigenvector
$v+U\in V/U\setminus\{0\}$. From Exercise~5B25~(a), we know that the minimal
polynomial of $T$ is a multiple of the minimal polynomial of $T/U$. Since
$\lambda$ is an eigenvalue of $T/U$, it must be a root of the minimal
polynomial of $T/U$. Therefore, $\lambda$ is also a root of the minimal
polynomial of $T$, which implies that $\lambda$ is an eigenvalue of $T$.
\end{solution}

\begin{exercise}
Suppose $\F=\R$, $V$ is finite-dimensional, and $T\in\Lin(V)$. Prove that the
minimal polynomial of $T_{\C}$ equals the minimal polynomial of $T$.
\exnote{The complexification $T_{\C}$ was defined in Exercise~33 of
Section~3B.}
\end{exercise}
\begin{solution}
Let $p(z)$ be the minimal polynomial of $T$. Then we have
\[ p(T)=0. \]
Consider the complexification $T_{\C}$ of $T$. By definition, $T_{\C}$
applied on $v+iw\in V_{\C}$ is given by
\[ T_{\C}(v+iw)=T(v)+iT(w) \]
for all $v,w\in V$. By induction, we can extend the polynomial $p(z)$ to act
on $V_{\C}$ by defining
\[ p(T_{\C})(v+iw)=p(T)v+ip(T)w. \]
Since $p(T)=0$, it follows that
\[ p(T_{\C})(v+iw)=0+i0=0 \]
for all $v,w\in V$. Therefore, $p(T_{\C})=0$, which shows that the minimal
polynomial of $T_{\C}$ divides $p(z)$.

Conversely, let $q(z)$ be the minimal polynomial of $T_{\C}$. We can split
$q(z)$ into its real and imaginary parts, say $q(z)=r(z)+is(z)$, where
$r(z),s(z)\in\Poly(\R)$. Then for any $v\in V$ which can be viewed as
$v+i0\in V_{\C}$, we have
\[ q(T_{\C})(v+i0)=r(T)(v+i0)+is(T)(v+i0)=r(T)v+is(T)v. \]
Since $q(T_{\C})=0$, it follows that $r(T)v=0$ and $s(T)v=0$ for all $v\in V$.
Therefore, both $r(T)$ and $s(T)$ must be $0$, which implies that the minimal
polynomial of $T$ divides $r(z)$ and $s(z)$. Since $q(z)=r(z)+is(z)$, it
follows that the minimal polynomial of $T$ divides $q(z)$. Combining both
directions, we conclude that the minimal polynomial of $T_{\C}$ equals the
minimal polynomial of $T$.
\end{solution}

\begin{exercise}
Suppose $V$ is finite-dimensional and $T\in\Lin(V)$. Prove that the minimal
polynomial of $T'\in\Lin(V')$ equals the minimal polynomial of $T$.
\exnote{The dual map $T'$ was defined in Section~3F.}
\end{exercise}
\begin{solution}
Let $p(z)$ be the minimal polynomial of $T$. Then we have
\[ p(T)=0. \]
Consider the dual map $T'\in\Lin(V')$ defined by $T'(\varphi)=\varphi\circ T$
for all $\varphi\in V'$. By induction, we can extend the polynomial $p(z)$ to
act on $V'$ by defining
\[ p(T')(\varphi)=p(T')(\varphi)=\varphi\circ p(T). \]
Since $p(T)=0$, it follows that
\[ p(T')(\varphi)=\varphi\circ0=0 \]
for all $\varphi\in V'$. Therefore, $p(T')=0$, which shows that the minimal
polynomial of $T'$ divides $p(z)$.

Conversely, let $q(z)$ be the minimal polynomial of $T'$. Then we have
\[ q(T')=0. \]
This implies that $q(T')(\varphi)=\varphi\circ q(T)=0$ for all
$\varphi\in V'$. By Exercise~3F32, we know that $V$ is isomorphic to $V''$,
the double dual of $V$. Using the isomorphism $\Lambda\colon V\to V''$ defined
in Exercise~3F32~(c), we can see that if $q(T')(\varphi)=0$ for all
$\varphi\in V'$, then
\[ \Lambda\bigl(q(T)v\bigr)(\varphi)=\varphi\bigl(q(T)v\bigr)
   =q(T')(\varphi)(v)=0 \]
for all $v\in V$ and $\varphi\in V'$. It follows that $q(T)v=0$ for all
$v\in V$. Therefore, $q(T)=0$, which shows that the minimal polynomial of $T$
divides $q(z)$. Combining both directions, we conclude that the minimal
polynomial of $T'$ equals the minimal polynomial of $T$.
\end{solution}

\begin{exercise}
Show that every operator on a finite-dimensional vector space of dimension at
least two has an invariant subspace of dimension two.
\exnote{Exercise~6 in Section~5C will give an improvement of this result when
$\F=\C$.}
\end{exercise}
\begin{solution}
Let $V$ be a finite-dimensional vector space with $\dim V\ge2$, and let
$T\in\Lin(V)$. Let $p(z)$ be the minimal polynomial of $T$. Over $\F$, every
nonconstant polynomial factors into linear and irreducible quadratic factors.
We split the proof into two cases based on the factorization of $p(z)$.
\begin{enumerate}[label=\arabic*)]
\item Suppose $p(z)$ has a linear factor $z-\lambda$ for some $\lambda\in\F$.
  If $\deg p(z)=1$, then $p(z)=z-\lambda$, which implies $T-\lambda I=0$, or
  simply $T=\lambda I$. Since $T$ is just a scalar multiple of the identity,
  every single subspace of $V$ is invariant under $T$. Because $\dim V\ge2$,
  we can choose any two linearly independent vectors $v,w\in V$, and
  $\spn(v,w)$ will form a $2$-dimensional invariant subspace.

  If $\deg p(z)\ge2$, then $p(z)/(z-\lambda)$ has degree at least $1$ and must
  itself contain either a linear factor or an irreducible quadratic factor.
  The latter is handled by Case~2 below, so we consider the two linear cases:
  \begin{itemize}
  \item Repeated linear factor: $p(z)=(z-\lambda)^2q(z)$. Since $p(z)$ is the
    minimal polynomial, there exists $v\in V$ such that
    $(T-\lambda I)q(T)v\ne0$ but $(T-\lambda I)^2q(T)v=0$. Let $w=q(T)v$ and
    $u=(T-\lambda I)w$. Then $u\ne0$ and
    $(T-\lambda I)u=(T-\lambda I)^2w=0$, which implies $Tu=\lambda u$. Thus,
    $u$ is an eigenvector corresponding to the eigenvalue $\lambda$.
    Moreover, $Tw=(T-\lambda I)w+\lambda w=u+\lambda w$, which shows that
    $Tw\in\spn(u,w)$.

    To confirm linear independence, suppose $u=\alpha w$ for some
    $\alpha\in\F$. Then $(T-\lambda I)w=\alpha w$, so $w$ is an eigenvector of
    $T$ with eigenvalue $\lambda+\alpha$. But then
    $(T-\lambda I)^2w=\alpha^2w$. Since $(T-\lambda I)^2w=0$ and $w\ne0$, this
    forces $\alpha=0$, contradicting our deduction that $\alpha\ne0$ (since
    $u\ne0$). Therefore, $u$ and $w$ are linearly independent, and
    $\spn(u,w)$ is a $2$-dimensional invariant subspace of $V$ under $T$.
  \item Distinct linear factor: $p(z)=(z-\lambda)(z-\mu)g(z)$ with
    $\lambda\ne\mu$. Then $T$ has two distinct eigenvalues, yielding two
    linearly independent eigenvectors $v$ and $u$. The subspace $\spn(v,u)$
    is $2$-dimensional, and since $Tv=\lambda v\in\spn(v,u)$ and
    $Tu=\mu u\in\spn(v,u)$, it is explicitly invariant.
  \end{itemize}
\item Suppose $p(z)$ has an irreducible quadratic factor, say $z^2+az+b$ for
  some $a,b\in\F$. This means we can write $p(z)=(z^2+az+b)q(z)$ for some
  polynomial $q(z)$. Since $p(z)$ is the minimal polynomial, $q(T)$ is not the
  zero operator, meaning there exists some vector $x\in V$ such that
  $q(T)x\ne0$.

  Let $u_1=q(T)x$. Evaluating the minimal polynomial identity gives:
  \[ 0=p(T)x=(T^2+aT+bI)q(T)x=(T^2+aT+bI)u_1. \]
  Rearranging this provides the operator relation $T^2u_1=-bu_1-au_2$, where
  we define $u_2=Tu_1$.

  The vectors $u_1$ and $u_2$ must be linearly independent. If $Tu_1=cu_1$ for
  some $c\in\F$, then $(T^2+aT+bI)u_1=(c^2+ac+b)u_1=0$. Since $u_1\ne0$, this
  implies $c^2+ac+b=0$, meaning $c$ is a root of $z^2+az+b$ over $\F$,
  contradicting its irreducibility.

  Let $U=\spn(u_1,u_2)$, which has dimension two. Checking the images of our
  basis elements under $T$ yields:
  \begin{align*}
  T(u_1) &= u_2\in U\\
  T(u_2) &= T^2u_1=-bu_1-au_2\in U.
  \end{align*}
  Since $T$ maps the generating elements of $U$ back into $U$, the subspace
  $U$ is a $2$-dimensional invariant subspace of $V$ under $T$.
\end{enumerate}
\end{solution}
\end{exc}
